\documentclass[10pt,a4paper]{amsart}
\usepackage[margin=19mm]{geometry}
\usepackage{amsmath,amssymb,mathtools}
\usepackage[scr=rsfs,cal=euler]{mathalfa}
\usepackage{xfrac}
\usepackage[unicode,hidelinks,bookmarks=false]{hyperref}
\newcommand{\ie}{i.e.\ }
\newcommand{\cf}{cf.\ }
\newcommand{\resp}{respectively\ }
\newcommand{\Subsection}{\subsection}
\providecommand{\nobreakdash}{\nobreak\hbox{-}}
\setlength{\emergencystretch}{2em}
\multlinegap0pt
\usepackage{amssymb}
\newtheorem{theorem}{Theorem}
\newtheorem{lemma}{Lemma}
\newtheorem{observation}{Observation}
\def\conv{\mbox{\rm conv}}
\title{On the maximum number of $k$-holes in point sets with no $(k + 1)$-hole}
\author{Andrew Suk, Su Zhou}
\date{Source: arXiv:2606.05721v1 (CC BY 4.0)}
\begin{document}
\maketitle

\noindent\emph{Source and licence.} On the maximum number of $k$-holes in point sets with no $(k + 1)$-hole. Version arXiv:2606.05721v1 (CC BY 4.0), distributed by arXiv under the Creative Commons Attribution 4.0 International licence, http://creativecommons.org/licenses/by/4.0/, as recorded in the arXiv licence metadata for this exact version and shown on the abstract page https://arxiv.org/abs/2606.05721v1. Full licence notice: \texttt{licenses/arxiv-2606.05721-CC-BY-4.0.txt}.\par\smallskip

\noindent\textit{Editorial scope.} Counted unit: the article's theorem hole1 (source0.tex lines 119--126). The article's complete original proof of it is delivered with it (source0.tex lines 230--418, one \texttt{proof} environment of 189 lines, which the article places away from the statement). No mathematics is written, repaired or re-derived: the statement and the proof are the article's own lines, in the article's own notation, and no retained line is substituted or re-flowed.  Retained uncounted as the source-local context the proof needs: lines 204--206.  One declared adaptation: exactly 3 ancillary strings inside the retained spans are omitted (image-inclusion commands and, where the source repeats a label, the repeated label definitions) and no other line differs from the source; the figure environments, with their placement, captions and labels, are kept, and the package records the omitted strings in fidelity receipts bound to the affected bodies.  No retained line contains a citation, so the wrapper carries no bibliography and no bibliography entry.  The retained lines contain the article's own diagram code, which the wrapper typesets with the standard tikz packages; no image file is delivered and the package declares no asset dependency.  Editorial formatting only, in addition: an AMS \texttt{amsart} preamble that restates the article's own declarations and macros used by the retained lines, namely the environments \texttt{theorem}, \texttt{lemma}, \texttt{observation} on separate counters, together with the standard packages those lines need.  The retained environments print in this document as Theorem 1, Lemma 1, Observation 1 in order of appearance, whereas the article prints the same items with the numbering of its own full document.  The complete native source of the article as distributed in the arXiv e-print for this exact version is bundled unchanged as \texttt{sources/202136/source0.tex}.
\begin{lemma}\label{lem:3cap}
Every Horton set of size $n$ has at least $\left\lfloor \frac{n-2}{2}\right\rfloor$ open 3-caps, and at least $\left\lfloor \frac{n-2}{2}\right\rfloor$ open 3-cups.
\end{lemma}

\begin{theorem}\label{hole1}
There are absolute constants $c_1,c_2>0$ such that the following holds. For every integer \(n>k\ge 6\),
\[
\left(\frac{c_1}{k}\right)^{\lfloor k/3\rfloor} n^{\lfloor k/3\rfloor}
\le h_k(n)\le
\left(\frac{c_2}{k}\right)^{\lceil k/2\rceil} n^{\lceil k/2\rceil}.
\]
\end{theorem}

\begin{proof}[Proof of Theorem \ref{hole1}] We start by proving the lower bound for $h_k(n)$.   Set \(m=\lfloor k/3\rfloor\), and consider the convex polygon \(C\) with \(m+1\) vertices forming a cap. More specifically, let the vertices of \(C\) be
\[
\left(\cos\left(\frac{\pi i}{m}\right),\sin\left(\frac{\pi i}{m}\right)\right)\in\mathbb{R}^2,
\]
for \(i=0,\ldots,m\). Thus, the upper boundary of \(C\) consists of \(m\) sides \(s_1,\ldots,s_m\), ordered from left to right.  Along each side \(s_i\), near its midpoint, we place a sufficiently small and flat copy \(H_i\) of a Horton set as follows.  Since there exist integers \(q\ge 0\) and \(0\le r<m\) such that
\[
n-(k\bmod 3)=q m+r
\]
where $(k\bmod 3)$ denotes the remainder of $k$ when divided by $3$, we let \(H_1,\dots,H_r\) be Horton sets of size \(q+1\), and \(H_{r+1},\dots,H_{m}\) be Horton sets of size \(q\). Then
\[
\sum_{i=1}^{m}|H_i|
=
r(q+1)+\left(m-r\right)q
=
qm+r
=
n-(k\bmod 3).
\]
Each such $H_i$ will be chosen so that \(H_i\) is a Horton set with respect to \(s_i\). Moreover, the points of each \(H_i\) are ordered by increasing \(x\)-coordinate, and this ordering is preserved after rotating the plane so that \(s_i\) becomes horizontal.




We set $P =  H_1\cup \cdots \cup H_m$, which lies above the $x$-axis. Let $P_0$ be a set of $k\bmod 3$ points placed on the $x$-axis very close to the origin and set $P^{\ast} = P_0\cup P$. By applying small perturbations, we may assume that $P^{\ast}$ is in general position. Hence $|P^{\ast}|=n$, and by choosing each copy of $H_i$ sufficiently small and flat, the set $P^{\ast}$ has the following properties.




\begin{enumerate}
    \item If $\ell$ is the line generated by any two points in $H_i$, then $P^{\ast}\setminus H_i$ lies below $\ell$.  

    
\item For each \(H_i=\{h_1,\ldots,h_{|H_i|}\}\), every line \(\ell\) passing through \(h_j\) and any other point of \(P^{\ast}\setminus H_i\) has the property that the points \(\{h_1,\ldots,h_{j-1}\}\) lie on one side of \(\ell\), while the points \(\{h_{j+1},\ldots,h_{|H_i|}\}\) lie on the other side.
\end{enumerate}





    \begin{figure}
        \centering
        
        \caption{Example with $k = 13$, $m = 4$ and $|P_0| = 1$.}
        \label{fig:const}
    \end{figure}

    

\noindent See Figure \ref{fig:const}.   We now have the following simple observation.



\begin{observation}\label{clean}
    Let $X \subset H_i$ be a cap with respect to segment $s_i$ and $p \in P^{\ast}\setminus H_i$.  Then $X\cup p$ is a hole in $P^{\ast}$ if and only if $X$ is an open cap with respect to $s_i$.  
\end{observation}




\begin{proof}
Let $X=\{h_{i_1},\ldots,h_{i_r}\}\subset H_i$ be ordered with respect to their \(x\)-coordinate and let $p \in P^{\ast}\setminus H_i$.  Since each $H_j$ is sufficiently small,  every point of \(P^{\ast}\setminus (H_i\cup \{p\})\) lies outside \(\conv(X\cup\{p\})\), so any point of \(P^{\ast}\) contained in the interior of \(\conv(X\cup\{p\})\) must belong to \(H_i\setminus X\).

 

Suppose first that \(X\) is an open cap with respect to \(s_i\). Then no point of \(H_i\setminus X\) lies in \(\conv(X\cup\{p\})\). Indeed, let \(h_t\in H_i\setminus X\). If \(t<i_1\) or \(t>i_r\), then property~(2) implies that \(h_t\) does not lie between the lines generated by \(ph_{i_1}\) and \(ph_{i_r}\).  Hence, \(h_t\notin \conv(X\cup\{p\})\). If \(i_1<t<i_r\), then, since \(X\) is an open cap, the point \(h_t\) lies above the polygonal chain \(h_{i_1}h_{i_2}\cdots h_{i_r}\). Thus again \(h_t\notin \conv(X\cup\{p\})\). It follows that \(\conv(X\cup\{p\})\) contains no point of \(P^{\ast}\setminus (X\cup\{p\})\), and hence \(X\cup\{p\}\) is a hole in \(P^{\ast}\).

Conversely, suppose that \(X\cup\{p\}\) is a hole in \(P^{\ast}\). If \(X\) were not an open cap, then there would exist a point
\(h_t\in H_i\setminus X\) with \(i_1<t<i_r\) lying below the polygonal chain
\(h_{i_1}h_{i_2}\cdots h_{i_r}\).
Since \(i_1<t<i_r\), property~(2) implies that \(h_t\) lies between the two
lines \(ph_{i_1}\) and \(ph_{i_r}\).
Hence \(h_t\in \conv(X\cup\{p\})\), a contradiction.\end{proof}















By Observation \ref{clean}, we can select an open 3-cap $X_i$ in each $H_i$, and together with $P_0$, to create a $k$-hole in $P^{\ast}$.  Indeed, by property (1), (2), $X_1\cup \cdots \cup X_m$ forms a cap, and together with $P_0$, we have $k$ points in convex position.  By Observation \ref{clean} and property (2), $X_1\cup \cdots \cup X_m \cup P_0$ is a hole.  Since each $H_i$ has $\Omega(n/m)$ open 3-caps by Lemma \ref{lem:3cap}, we obtain at least $\Omega((n/m)^{\lfloor k/3\rfloor})$ holes of size $k$ in $P^{\ast}$.  Since $ m = \lfloor k/3\rfloor$, there is an absolute constant $c_1 > 0$ such that   

$$h_k(n) \geq \left(\frac{c_1}{k}\right)^{\lfloor k/3\rfloor }n^{\lfloor k/3\rfloor}.$$ 


Moreover, \(P^{\ast}\) does not contain a \((k+1)\)-hole.  Indeed, suppose
for contradiction that \(X\subset P^{\ast}\) is a \((k+1)\)-hole.  Since each
\(H_i\) contains no \(7\)-hole, the set \(X\) is not contained in any single
\(H_i\).  Moreover, if \(|X\cap H_i|\le 3\) for every \(i\), then
\[
        |X|\le 3m+|P_0|=k,
\]
a contradiction.  Hence, for some \(i\), we have \(|X\cap H_i|\ge 4\) and
\(X\cap(P^{\ast}\setminus H_i)\neq\emptyset\).  Thus \(X\cap H_i\) must form
a cap with respect to \(s_i\).  Since \(H_i\) is lower \(4\)-closed with
respect to \(s_i\), Observation~\ref{clean} implies that \(X\) is not a
\((k+1)\)-hole, a contradiction.


 
\medskip

Next, we prove the upper bound for \(h_k(n)\). Choose \(c_2>0\) sufficiently large so that \((c_2/2e)^3>2\).  For the sake of contradiction, suppose \(P\) is a point set of size \(n'\le n\) in the plane with no \((k+1)\)-hole, but with at least \(\left(c_2/k\right)^{\lceil k/2\rceil} n^{\lceil k/2\rceil}\) \(k\)-holes, where \(c_2\) is a sufficiently large constant. Since \(n'\le n\), the set \(P\) has at least \(\left(c_2/k\right)^{\lceil k/2\rceil}(n')^{\lceil k/2\rceil}\) \(k\)-holes. After rotation, we may assume that no two points from $P$ share the same $x$-coordinate. 

Let $X = \{x_1,\ldots, x_k\}\subset P$ be a convex $k$-gon in $P$ such that $x_1$ is the left-most point in $X$, and the elements in $X$ are ordered in clockwise order.  We say that $Y\subset X$ is a \emph{support} of $X$ if $Y  = \{x_1,x_3,\ldots, x_{k-1}\}$ when $k$ is even, and $Y = \{x_1,x_3,\ldots, x_k\}$ when $k$ is odd. Hence, each $k$-hole $X$ in $P$ corresponds to a unique support $Y\subset X$.  Since $|Y| = \lceil k/2\rceil$, and since \(P\) contains at least \(\left(c_2/k\right)^{\lceil k/2\rceil}(n')^{\lceil k/2\rceil}\) \(k\)-holes, there must be a subset $Y\subset P$ of size $\lceil k/2\rceil$ such that $Y$ is the support for at least

\begin{align*}
\frac{\left(c_2/k\right)^{\lceil k/2\rceil} (n')^{\lceil k/2\rceil}}
{\binom{n'}{\lceil k/2\rceil}}
&\ge
\left(\frac{c_2}{k}\right)^{\lceil k/2\rceil}\lceil k/2\rceil!\, 
\\
&>
\left(\frac{c_2}{k}\right)^{\lceil k/2\rceil}
\left(\frac{\lceil k/2\rceil}{e}\right)^{\lceil k/2\rceil}  \\
&=
\left(\frac{c_2\lceil k/2\rceil}{ek}\right)^{\lceil k/2\rceil}  \\
&\ge
\left(\frac{c_2}{2e}\right)^{\lceil k/2\rceil} 
>2,
\end{align*}

\noindent $k$-holes in $P$.  
Let $X,X' \subset P$ be two distinct $k$-holes in $P$ that use $Y$ as a support.  Hence, we have $X = \{x_1,\ldots, x_k\}$, $X' = \{x'_1,\ldots, x'_k\}$, $x_1 = x'_1$, $x_3 = x'_3$, etc. Just as above, $x_1$ is the left-most point in $X$, and the elements in $X$ are labeled in clockwise order.  Likewise, \(x'_1=x_1\) is the left-most point of \(X'\), and the elements of \(X'\) are labeled in clockwise order. We are going to finish the proof using the following observation. 


\begin{observation}\label{diff}

Each $Y\subseteq P$ with $|Y|=\lceil k/2\rceil$ supports at most one $k$-hole $X\subseteq P$.  
    
\end{observation}

\begin{proof}

    \begin{figure}
        \centering
        
        \caption{Regions $T_i$ shaded in gray for $k = 6$, and $k = 7$.}
        \label{exconv}
    \end{figure}

Suppose $Y$ supports two distinct $k$-holes $X$ and $X'$ as above. We define regions \[T_1,\ldots,T_{\lfloor (k+1)/2\rfloor}\] outside of \(\conv(Y)\) as follows. If \(k\) is even, then for each \(1\le i\le k/2\), let \(T_i\) be the region outside of \(\conv(Y)\) bounded by the segment \(x_{2i-1}x_{2i+1}\) and by the two lines \(x_{2i-3}x_{2i-1}\) and \(x_{2i+1}x_{2i+3}\), where indices are taken modulo \(k\).  Now suppose that \(k\) is odd. For each \(1\le i\le (k-1)/2\), let \(T_i\) be the region outside of \(\conv(Y)\) bounded by the segment \(x_{2i-1}x_{2i+1}\) and by the two lines containing the adjacent sides of \(\conv(Y)\), and let \(T_{(k+1)/2}\) be the region outside of \(\conv(Y)\) bounded by the segment \(x_1x_k\) and by the two lines \(x_{k-2}x_k\) and \(x_1x_3\).  See Figure \ref{exconv}.





  First, let us consider the case when $k$ is even.  Suppose $X,X'$ share some point not in $Y$.  Let $i\leq k/2$ be the maximum integer such that $x_{2i} \neq x'_{2i}$ and $x_{2i + 2} = x'_{2i+2}$, where $x_{k + 1} = x_1$ and $x_{k + 2} = x_2$.  We define $S_i$ to be the region inside of $T_i$, bounded by the lines $x_{2i-3}x_{2i-1}$, $x_{2i-2}x_{2i-1}$, and $x_{2i}x_{2i+1}$.   Likewise, we define $R_i$ to be the region inside of $T_i$, bounded by the lines $x_{2i-1}x_{2i}$, $x_{2i+1}x_{2i+2}$, and $x_{2i + 1}x_{2i+3}$.  See Figure \ref{fig:es1}.    





We start in region $T_i$.  Since both $X$ and $X'$ are $k$-holes, and since $P$ does not contain a $(k + 1)$-hole, this implies that $x'_{2i}$ must lie in $S_i$ or $R_i$: otherwise, we must have that $X\cup \{x'_{2i}\}$ is a $(k+1)$-hole, or $\conv(\{x'_{2i - 1},x'_{2i},x'_{2i + 2}\})$/$\conv(\{x_{2i - 1},x_{2i},x_{2i + 2}\})$ is not empty.  See Figure \ref{fig:es1}.  If $x'_{2i} \in R_i$, then $x'_{2i + 1} \in \conv(\{x'_{2i - 1},x'_{2i},x'_{2i + 2}\})$ since $x_{2i + 2} = x'_{2i+2}$.  This contradicts the fact that $X'$ is a $k$-hole.  Therefore, we must have $x'_{2i} \in S_i$ and we proceed to region $T_{i-1}$.  

We apply the same argument inside region $T_{i-1}$.  Again, we must have $x'_{2i - 2} \neq x_{2i - 2}$ and $x'_{2i - 2}  \notin R_{i-1}$, since otherwise $x'_{2i - 1} \in \conv(\{x'_{2i - 3},x'_{2i-2},x'_{2i}\})$ and $X'$ would not be a $k$-hole.  Hence, we must have $x'_{2i-2} \in S_{i-1}$ and we proceed to region $T_{i-2}$.  Iterating this argument in counterclockwise order over all indices forces $x_{2i+2} \neq x'_{2i+2}$, contradicting the choice of $i$. Hence, we must have that $x'_{2i}\in S_i$ or $x'_{2i}\in R_i$ for every $i$. However, then $x'_2$ or $x'_{k}$ will be the left-most vertex of $X'$, which contradicts the fact that $x_1$ is the left-most vertex of $X'$. 

\medskip


Now suppose that $k$ is odd, which implies that the interior of region $T_{\lceil k/2\rceil}$ does not contain a point from $X$ and $X'$. Since both $X$ and $X'$ are $k$-holes, and since $P$ does not contain a $(k + 1)$-hole, this implies that $x_{k-1}=x'_{k-1}$ or $x'_{k-1} \in S_{\lceil k/2\rceil-1}$. As $X\neq X'$, let $i\leq (k-1)/2$ be maximal such that $x_{2i} \neq x'_{2i}$. 





We start in region $T_i$, and argue just as above.  Since both $X$ and $X'$ are $k$-holes, and since $P$ does not contain a $(k + 1)$-hole, this implies that $x'_{2i}$ must lie in $S_i$ or $R_i$.  However, since we are assuming $i$ is maximal such that $x_{2i} \neq x'_{2i}$, this implies that $x'_{2i} \notin R_i$ since otherwise $x'_{2i + 1} \in \conv(\{x'_{2i - 1},x'_{2i},x'_{2i + 2}\})$ and $X'$ would not be a $k$-hole.  Therefore, $x'_{2i} \in S_i$ and we proceed to region $T_{i - 1}$.   

  

    \begin{figure}
        \centering
        
        \caption{Regions $S_i$ and $R_i$ shaded in gray.}
        \label{fig:es1}
    \end{figure}

If $i - 1 \geq 2$, then again, we must have $x'_{2i - 2} \neq x_{2i - 2}$ since otherwise $$x'_{2i - 1} \in \conv(\{x'_{2i - 3},x'_{2i-2},x'_{2i}\})$$ and $X'$ would not be a $k$-hole.  By the same arguments as above, this implies that $x'_{2i - 2}$ must lie in $S_{i - 1}$ and we proceed to region $T_{i - 2}$.  Iterating this argument cyclically in counterclockwise order leads us to region $T_1$ with $x_2 \neq x'_2$ and $x'_2 \not\in R_1$ and $S_1  = \emptyset$.  Since both $X$ and $X'$ are $k$-holes, and neither uses a point in $T_{\lceil k/2\rceil}$, $X\cup \{x'_2\}$ is a $(k + 1)$-hole.  Contradiction.\end{proof}



This completes the proof.\end{proof}

\end{document}
